\documentclass[12pt,a4paper]{article}

\usepackage[utf8]{inputenc}
\usepackage[T1]{fontenc}
\usepackage[italian]{babel}
\usepackage{geometry}
\usepackage{amsmath}
\usepackage{amssymb}
\usepackage{booktabs}
\usepackage{enumitem}
\usepackage{xcolor}
\usepackage{titlesec}
\usepackage{fancyhdr}
\usepackage{mdframed}
\usepackage{hyperref}
\usepackage{bookmark}
\usepackage{parskip}
\usepackage{array}
\usepackage{listings}
\usepackage{textcomp}

\geometry{top=2.5cm, bottom=2.5cm, left=3cm, right=3cm, headheight=14.5pt}

% Colori
\definecolor{primary}{RGB}{0, 90, 160}
\definecolor{teal}{RGB}{0, 128, 128}
\definecolor{darkgray}{RGB}{70, 70, 70}
\definecolor{lightgray}{RGB}{245, 245, 245}
\definecolor{codebg}{RGB}{248, 248, 248}
\definecolor{codeframe}{RGB}{180, 200, 220}
\definecolor{green}{RGB}{0, 120, 60}

% Titoli sezione
\titleformat{\section}{\large\bfseries\color{primary}}{}{0em}{\thesection.~}[\titlerule]
\titleformat{\subsection}{\normalsize\bfseries\color{teal}}{}{0em}{}
\titleformat{\subsubsection}{\normalsize\itshape\color{darkgray}}{}{0em}{}

% Header/Footer
\pagestyle{fancy}
\fancyhf{}
\fancyhead[L]{\small\color{darkgray}\textit{TLN -- Tecnologie del Linguaggio Naturale}}
\fancyhead[R]{\small\color{darkgray}\textit{A.A. 2025/26 -- Parte 2 (Radicioni)}}
\fancyfoot[C]{\thepage}
\renewcommand{\headrulewidth}{0.4pt}

% Box definizione
\newmdenv[backgroundcolor=lightgray,linecolor=primary,linewidth=1.5pt,
  innerleftmargin=10pt,innerrightmargin=10pt,innertopmargin=8pt,
  innerbottommargin=8pt,skipabove=6pt,skipbelow=6pt]{defbox}

% Box formula
\newmdenv[backgroundcolor=white,linecolor=teal,linewidth=1pt,
  leftmargin=10pt,rightmargin=10pt,innerleftmargin=10pt,
  innerrightmargin=10pt,innertopmargin=6pt,innerbottommargin=6pt,
  skipabove=4pt,skipbelow=4pt]{formulabox}

% Box codice
\lstset{
  backgroundcolor=\color{codebg},
  frame=single, rulecolor=\color{codeframe},
  basicstyle=\small\ttfamily,
  keywordstyle=\color{primary}\bfseries,
  commentstyle=\color{green}\itshape,
  stringstyle=\color{teal},
  language=Python,
  breaklines=true, breakatwhitespace=true,
  showstringspaces=false,
  numbers=none,
  xleftmargin=8pt, xrightmargin=8pt,
  aboveskip=6pt, belowskip=6pt
}

\hypersetup{colorlinks=true,linkcolor=primary,urlcolor=teal,
  pdftitle={TLN - Parte 2 Radicioni A.A. 2025/26}}

% =====================================================================
\begin{document}

\begin{titlepage}
  \centering
  \vspace*{2cm}
  {\color{primary}\rule{\linewidth}{2pt}}\\[0.5cm]
  {\Huge\bfseries\color{primary} Tecnologie del Linguaggio Naturale}\\[0.4cm]
  {\Large\bfseries\color{darkgray} Riassunto -- Parte 2 (Radicioni)}\\[0.3cm]
  {\color{primary}\rule{\linewidth}{2pt}}\\[1cm]
  {\large Università degli Studi di Torino (UniTO)}\\[0.2cm]
  {\large Corso di Laurea Magistrale in Informatica}\\[0.2cm]
  {\large A.A. 2025/26}\\[2cm]
  {\large\itshape Alessandro Olivero}\\[0.5cm]
  {\normalsize\color{darkgray} Argomenti: KR, WordNet, WSD, N-grammi, VSM, Word2Vec, Transformers}\\[3cm]
  {\small\color{darkgray} Documento generato a scopo di studio}
  \vfill
\end{titlepage}

\tableofcontents
\newpage

% =====================================================================
\section{Introduzione alla Semantica Lessicale e al NLP}

La \textbf{semantica lessicale} studia il significato delle parole, le relazioni tra di esse, e come rappresentare e computare questi significati.

\subsection{Le quattro generazioni dei modelli linguistici}

\begin{center}
\begin{tabular}{>{\bfseries}llll}
\toprule
Era & Tipo & Approccio & Metodi chiave \\
\midrule
anni '90 & Statistical LM & Assunzione di Markov & N-grammi, smoothing \\
2013 & Neural LM & Reti neurali & Word2Vec (rappresentazioni distribuite) \\
2018 & Pre-trained LM & Transformer & BERT, GPT (bidirezionale/unidirezionale) \\
2020+ & Large LM & Scaling + prompting & GPT-3/4, ChatGPT, Claude \\
\bottomrule
\end{tabular}
\end{center}

\subsection{Livelli di analisi linguistica}

\begin{enumerate}[leftmargin=1.5em, label=\arabic*.]
  \item \textbf{Fonetica/fonologia}: suoni e loro combinazioni.
  \item \textbf{Morfologia}: struttura interna delle parole (morfemi, flessione, derivazione).
  \item \textbf{Sintassi}: struttura delle frasi (costituenti, dipendenze).
  \item \textbf{Semantica}: significato delle parole e delle frasi.
  \item \textbf{Pragmatica}: uso del linguaggio nel contesto (speech acts, implicature).
  \item \textbf{Discorso}: struttura testuale, coerenza, coreference.
\end{enumerate}

\subsection{Ambiguità nel linguaggio naturale}

\begin{itemize}[leftmargin=1.5em]
  \item \textbf{Morfologica/PoS}: ``back'' = ADJ / NOUN / VERB / ADV.
  \item \textbf{Sintattica}: PP attachment, coordinazione (vedi Parte 1).
  \item \textbf{Semantica}: polisemia (``banca''), scope quantificazionale.
  \item \textbf{Fonetica}: omofoni.
\end{itemize}

\textbf{Esempio classico (Mazzei/Radicioni)}: ``I made her duck.''
\begin{itemize}[leftmargin=1.5em, nosep]
  \item \emph{duck} = sostantivo (anatra) oppure verbo (abbassarsi)
  \item \emph{make} = causativo (le ho fatto abbassare) oppure creativo (le ho cucinato l'anatra)
\end{itemize}

% =====================================================================
\section{Knowledge Representation (KR) e Ontologie}

\subsection{Ontologia: definizione}

Il termine \emph{ontologia} viene dal greco: \textit{on} (essere) + \textit{logos} (discorso). In filosofia studia \emph{cosa esiste} e la struttura della realtà. In informatica:

\begin{defbox}
\textbf{Ontologia} (informatica): specificazione esplicita e formale di una \emph{concettualizzazione}.\\[4pt]
\textbf{Concettualizzazione}: struttura formale di una porzione della realtà così come è percepita e organizzata da un agente, \emph{indipendentemente} dal vocabolario usato e dall'occorrenza di situazioni specifiche.
\end{defbox}

\subsection{Livelli di precisione ontologica}

In ordine crescente di formalizzazione:

\begin{enumerate}[leftmargin=1.5em, label=\arabic*.]
  \item \textbf{Catalogo}: lista di termini (minima precisione).
  \item \textbf{Glossario}: lista con definizioni informali.
  \item \textbf{Thesaurus}: glossario + relazioni lessicali (sinonimia, antonimia).
  \item \textbf{Tassonomia}: classificazione gerarchica (relazioni \emph{is-a}).
  \item \textbf{Schema DB / OO}: struttura formale con proprietà e vincoli.
  \item \textbf{Teoria assiomatica}: specificazione logica completa (massima precisione).
\end{enumerate}

\textbf{Esempio -- dominio Tennis}: catalogo $\to$ \{tennis, campo, rete\} $\to$ tassonomia \texttt{gioco\_atletico} $\to$ \texttt{gioco\_di\_campo} $\to$ \texttt{tennis} $\to$ teoria: $\mathit{tennis}(x) \Rightarrow \mathit{gioco\_campo}(x) \wedge \exists y[\mathit{campo}(y) \wedge \mathit{giocato\_in}(x,y)]$.

\subsection{Componenti di un'ontologia}

\begin{itemize}[leftmargin=1.5em]
  \item \textbf{Componente terminologica}: classi, relazioni, assiomi -- indipendente da situazioni specifiche.
  \item \textbf{Componente assertiva}: fatti individuali -- riflette stati epistemici specifici.
\end{itemize}

\subsection{DOLCE -- Descriptive Ontology for Linguistic and Cognitive Engineering}

DOLCE è un'ontologia di alto livello (upper ontology) che fornisce un framework per descrivere entità del mondo reale.

\subsubsection{Distinzione fondamentale: Endurant vs Perdurant}

\begin{defbox}
\begin{itemize}[leftmargin=1.5em]
  \item \textbf{Endurant (Oggetto)}: è \emph{completamente presente} in qualsiasi istante in cui esiste. Può cambiare nel tempo mantenendo la propria identità. Es.: una rosa.
  \item \textbf{Perdurant (Evento)}: non è completamente presente in ogni istante; si \emph{svolge} nel tempo e non cambia (tutte le sue parti sono essenziali). Es.: una conferenza, una performance musicale.
\end{itemize}
\end{defbox}

\textbf{Relazione}: gli Endurant \emph{partecipano} nei Perdurant.

\subsubsection{Identità e criteri di identità}

Un criterio di identità determina cosa rende unica un'entità. Esempio: \texttt{time-duration} è identificata dalla lunghezza; \texttt{time-interval} è identificata da inizio e fine. Pur avendo parti comuni, \emph{non} sono in relazione is-a perché hanno criteri di identità diversi.

\subsubsection{Qualità, Quale e Quality Space}

\begin{itemize}[leftmargin=1.5em]
  \item \textbf{Quality (particolare)}: il colore di \emph{questa specifica} rosa.
  \item \textbf{Quale (valore)}: una specifica tonalità di rosso.
  \item \textbf{Quality Space (universale)}: l'insieme di tutti i valori di colore possibili.
  \item Relazione: la qualità \emph{inheres in} (dipende da) l'entità.
\end{itemize}

\subsubsection{Tassonomia DOLCE (semplificata)}

\begin{center}
\begin{tabular}{ll}
\toprule
\textbf{Categoria} & \textbf{Sottocategorie} \\
\midrule
ENDURANT & Fisico (oggetti, materia) / Non-fisico (mentale, sociale) \\
PERDURANT & Statico (State, Process) / Dinamico (Achievement, Accomplishment) \\
QUALITY & Fisica (spaziale, \ldots) / Temporale \\
REGION & Quality region, Space region, Time region \\
\bottomrule
\end{tabular}
\end{center}

\subsubsection{La relazione part-of}

La relazione di \emph{parthood} ha diverse varianti:
\begin{itemize}[leftmargin=1.5em, nosep]
  \item Componente (il manico è componente della porta)
  \item Ingrediente (la farina è ingrediente del pane)
  \item Porzione (una fetta di pane)
  \item Area (una città in una regione)
  \item Membro (una nave in una flotta)
  \item Pezzo (un manico staccato)
\end{itemize}

% =====================================================================
\section{WordNet: Struttura e Applicazioni}

\subsection{Cos'è WordNet}

WordNet (Miller et al., 1990--) è un \textbf{database lessicale} dell'inglese che organizza nomi, verbi, aggettivi e avverbi per \emph{significato} anziché per forma.

\begin{defbox}
\textbf{Dimensioni}: $\sim$117.000 synset, $\sim$155.000 coppie parola--senso.
\end{defbox}

\subsection{Il concetto di Synset}

\begin{defbox}
Un \textbf{synset} (synonym set) è un insieme di parole che esprimono lo stesso concetto in un dato contesto. Ogni parola può appartenere a più synset (polisemia).
\end{defbox}

\textbf{Esempio}: \texttt{car.n.01} = \{car, auto, automobile, machine, motorcar\}; ognuno esprime lo stesso concetto di ``veicolo a motore a quattro ruote''.

\subsection{Relazioni lessicali e semantiche}

\begin{center}
\begin{tabular}{lll}
\toprule
\textbf{Relazione} & \textbf{Tipo} & \textbf{Esempio} \\
\midrule
Sinonimia & Lessicale & car = automobile \\
Antonimia & Lessicale & good $\neq$ bad \\
Iponimia & Semantica & car \textbf{is-a} vehicle \\
Iperonimia & Semantica & vehicle is-hypernym-of car \\
Meronimia & Semantica & wheel \textbf{is-part-of} car \\
Olonimia & Semantica & car \textbf{has-part} wheel \\
Entailment & Semantica & to buy entails to pay \\
Causazione & Semantica & to kill causes to die \\
Similar-to & (agg.) & happy \textit{similar-to} joyful \\
\bottomrule
\end{tabular}
\end{center}

\subsection{Misure di similarità su WordNet}

\subsubsection{Misure path-based}

\textbf{Shortest Path Similarity}:
\begin{formulabox}
\[
\text{sim}_{\text{path}}(w_1, w_2) = \frac{1}{\text{shortestPath}(w_1,w_2) + 1}
\]
\end{formulabox}
Limite: non considera la profondità della gerarchia (due nodi vicini in cima alla gerarchia sembrano molto simili).

\textbf{Leacock-Chodorow (LC)}:
\begin{formulabox}
\[
\text{sim}_{\text{LC}}(w_1, w_2) = -\log\!\left(\frac{\text{shortestPath}}{2 \times D}\right)
\]
\end{formulabox}
dove $D$ è la profondità massima della tassonomia.

\textbf{Wu-Palmer (WUP)} -- \emph{usata nel corso}:
\begin{formulabox}
\[
\text{sim}_{\text{WUP}}(w_1, w_2) = \frac{2 \times \text{depth}(\text{LCA})}{\text{depth}(w_1) + \text{depth}(w_2)}
\]
\end{formulabox}
dove LCA = Lowest Common Ancestor (antenato comune più basso). Range: $[0, 1]$. Premia le coppie con LCA profondo.

\textbf{Esempio Wu-Palmer}:

\begin{center}
\begin{tabular}{lllll}
\toprule
Coppia & $\text{depth}(w_1)$ & $\text{depth}(w_2)$ & $\text{depth}(\text{LCA})$ & $\text{sim}_{\text{WUP}}$ \\
\midrule
PUG -- POODLE & 3 & 3 & 2 (DOG) & $4/6 = 0.667$ \\
PUG -- SIAMESE & 3 & 3 & 1 (ANIMAL) & $2/6 = 0.333$ \\
\bottomrule
\end{tabular}
\end{center}

\subsubsection{Misure basate sull'Information Content (IC)}

L'\textbf{Information Content} di un concetto $c$ è:
\[
\text{IC}(c) = -\log P(c)
\]
dove $P(c)$ è la probabilità che una parola nel corpus sia iponimo di $c$.

\textbf{Resnik}:
\begin{formulabox}
\[
\text{sim}_{\text{Res}}(w_1, w_2) = -\log P(\text{LCA}) = \text{IC}(\text{LCA})
\]
\end{formulabox}
Considera solo il LCA; non dipende dalla distanza tra $w_1$, $w_2$ e il LCA.

\textbf{Lin}:
\begin{formulabox}
\[
\text{sim}_{\text{Lin}}(w_1, w_2) = \frac{2 \times \text{IC}(\text{LCA})}{\text{IC}(w_1) + \text{IC}(w_2)}
\]
\end{formulabox}
Versione normalizzata di Resnik. Range: $[0,1]$.

\textbf{Jiang-Conrath}:
\begin{formulabox}
\[
\text{sim}_{\text{JC}}(w_1, w_2) = \frac{1}{\text{IC}(w_1) + \text{IC}(w_2) - 2 \times \text{IC}(\text{LCA})}
\]
\end{formulabox}

% =====================================================================
\section{Word Sense Disambiguation (WSD)}

\subsection{Definizione del problema}

La \textbf{polisemia} è la proprietà per cui una parola ha più significati. Il WSD è il task di determinare \emph{quale} significato una parola ha in un dato contesto.

\begin{defbox}
\textbf{Esempio}: ``bank'' nel testo ``I deposited money in the bank'' $\to$ istituto finanziario (non riva del fiume).
\end{defbox}

\textbf{Baseline}: assegnare sempre il senso più frequente (Most Frequent Sense, MFS) $\to$ $\sim$65\% di accuratezza su inglese.

\subsection{Algoritmo di Lesk (1986)}

L'algoritmo di Lesk è un metodo \textbf{knowledge-based} che sfrutta le definizioni del dizionario.

\begin{defbox}
\textbf{Idea}: il senso corretto è quello la cui definizione (gloss) presenta la maggior \emph{sovrapposizione} con il contesto della parola da disambiguare.
\end{defbox}

\textbf{Algoritmo}:
\begin{enumerate}[leftmargin=1.5em, label=\arabic*.]
  \item Per ogni senso $s$ della parola $w$:
  \item Ottieni il gloss $G(s)$ dal dizionario (WordNet).
  \item Calcola l'overlap con il contesto $C$ della frase: $\text{score}(s) = |G(s) \cap C|$.
  \item Ritorna il senso con score massimo.
\end{enumerate}

\begin{formulabox}
\[
\hat{s} = \arg\max_{s} \bigl|\text{gloss}(s) \cap \text{context}(w)\bigr|
\]
\end{formulabox}

\textbf{Esempio}: ``bank'' in ``I deposited money in the bank'':
\begin{itemize}[leftmargin=1.5em, nosep]
  \item Gloss senso 1 (finanziario): ``financial institution for money and deposits''
  \item Gloss senso 2 (fiume): ``edge of a river''
  \item Contesto: \{deposited, money, bank\}
  \item Overlap: senso 1 >> senso 2 $\to$ scelto senso 1.
\end{itemize}

\textbf{Varianti}:
\begin{itemize}[leftmargin=1.5em, nosep]
  \item \textbf{Extended Lesk}: include definizioni delle parole correlate (iponimi, iperoni\-mi).
  \item \textbf{Normalized Lesk}: divisione per la lunghezza del gloss.
\end{itemize}

\subsection{Metodi graph-based}

Costruzione di un grafo dei sensi candidati:
\begin{itemize}[leftmargin=1.5em]
  \item Nodi = sensi candidati di tutte le parole nella frase.
  \item Archi = relazioni semantiche in WordNet (con pesi = similarità).
  \item Propagazione del contesto tramite algoritmi come \textbf{PageRank}.
\end{itemize}

\begin{formulabox}
\[
\text{rank}(s) = \frac{1-d}{N} + d \sum_{t \in \text{vicini}(s)} \frac{\text{rank}(t)}{\text{outdegree}(t)}
\]
\end{formulabox}
dove $d \approx 0.85$ è il damping factor.

\subsection{Metodi knowledge-based (combinazione di sensi)}

Per ogni frase, si cercano le assegnazioni di sensi $\{s_1,\dots,s_k\}$ che massimizzano la similarità complessiva tra tutte le coppie:

\[
\hat{S} = \arg\max_{\{s_1,\dots,s_k\}} \sum_{i \neq j} \text{sim}(s_i, s_j)
\]

\subsection{Metodi supervisionati}

Usano corpus annotati con i sensi corretti (SemCor, SemEval). Feature: parole vicine, PoS tag, collocazioni. Classificatori: SVM, Naive Bayes, reti neurali, BERT fine-tuned. Problema: bottleneck dell'annotazione (pochi dati per sensi rari).

\subsection{Sfide del WSD}

\begin{itemize}[leftmargin=1.5em]
  \item Sparsità dei dati (rari sensi poco osservati).
  \item Distinguere polisemia (stessa parola, sensi correlati) da omografia (stessa forma, origini diverse: ``bank'' finanziario vs geografico).
  \item Distinzioni fine-grained difficili (es. significati molto simili).
  \item Corpus di addestramento limitati per lingue non inglesi.
\end{itemize}

% =====================================================================
\section{Text Processing: Pipeline di Base}

\subsection{Tokenizzazione}

La \textbf{tokenizzazione} suddivide il testo in unità minime significative (token).

\begin{itemize}[leftmargin=1.5em]
  \item \textbf{Word tokenization}: separazione per spazi e punteggiatura.
  \item \textbf{Sentence tokenization}: separazione in frasi (`.', `!', `?').
  \item \textbf{Subword tokenization}: BPE, WordPiece (per Transformers).
\end{itemize}

\textbf{Sfide}: contrazioni (\textit{don't}), trattini (\textit{mother-in-law}), numeri (\textit{\$1,234.56}), abbreviazioni (\textit{Dr.}), emoji, social media (hashtag, mention).

\textbf{Implementazione NLTK}:
\begin{lstlisting}
from nltk.tokenize import word_tokenize, sent_tokenize
tokens = word_tokenize("Hello, world!")  # ['Hello', ',', 'world', '!']
frasi  = sent_tokenize("Prima. Seconda.") # ['Prima.', 'Seconda.']
\end{lstlisting}

\subsection{Normalizzazione}

\begin{itemize}[leftmargin=1.5em]
  \item \textbf{Lowercasing}: ``The'' $\to$ ``the''. Riduce la sparsità; perde info su propri.
  \item \textbf{Rimozione punteggiatura/caratteri speciali}: dipende dal task.
\end{itemize}

\subsection{Stemming}

Lo \textbf{stemming} riduce le parole alla radice (stem) applicando regole morfologiche.

\textbf{Porter Stemmer} (algoritmo più comune): applica regole in cascata per rimuovere suffissi.

\begin{center}
\begin{tabular}{ll}
\toprule
Input & Output (stem) \\
\midrule
inflations & inflat \\
consolingly & consol \\
computing & comput \\
\bottomrule
\end{tabular}
\end{center}

\textbf{Limiti}: over-stemming (``computer'' e ``computation'' $\to$ ``comput''); under-stemming; output non è una parola valida.

\subsection{Lemmatizzazione}

La \textbf{lemmatizzazione} riduce una parola al suo \emph{lemma} (forma canonica del dizionario), sfruttando informazioni morfologiche e il PoS tag.

\begin{center}
\begin{tabular}{lll}
\toprule
Input & PoS & Lemma \\
\midrule
running & VERB & run \\
better & ADJ & good \\
better & VERB & better \\
was & VERB & be \\
\bottomrule
\end{tabular}
\end{center}

\begin{lstlisting}
from nltk.stem import WordNetLemmatizer
lem = WordNetLemmatizer()
lem.lemmatize("running", pos="v")  # "run"
lem.lemmatize("better",  pos="a")  # "good"
\end{lstlisting}

\subsection{Stemming vs Lemmatizzazione}

\begin{center}
\begin{tabular}{lll}
\toprule
\textbf{Aspetto} & \textbf{Stemming} & \textbf{Lemmatizzazione} \\
\midrule
Input richiesto & Parola & Parola + PoS \\
Output & Stem (può non essere valido) & Lemma (parola valida) \\
Velocità & Veloce & Lento \\
Precisione & $\sim$70\% & $\sim$95\% \\
Uso tipico & Information Retrieval & NLP avanzato, analisi \\
\bottomrule
\end{tabular}
\end{center}

\subsection{Rimozione Stopword}

Le \textbf{stopword} sono parole molto frequenti con poco contenuto semantico (``a'', ``the'', ``is'', ``and'', \ldots). La loro rimozione riduce il rumore e accelera il processing; tuttavia può eliminare informazioni importanti per certi task (es.\ ``not'').

\begin{lstlisting}
from nltk.corpus import stopwords
stop = set(stopwords.words('english'))
tokens_filtrati = [t for t in tokens if t not in stop]
\end{lstlisting}

% =====================================================================
\section{N-grammi e Modelli di Linguaggio}

\subsection{Definizioni}

\begin{defbox}
Un \textbf{n-gramma} è una sequenza di $n$ token consecutivi.\\
\textbf{Unigramma} ($n=1$): singola parola.\\
\textbf{Bigramma} ($n=2$): coppia di parole consecutive.\\
\textbf{Trigramma} ($n=3$): tripla di parole consecutive.
\end{defbox}

\subsection{Assunzione di Markov}

\textbf{Chain rule} esatta:
\[
P(w_1, w_2, \dots, w_n) = P(w_1)\cdot P(w_2|w_1)\cdot P(w_3|w_1,w_2)\cdots P(w_n|w_1,\dots,w_{n-1})
\]

\textbf{Approssimazione di Markov} (per bigrammi):
\begin{formulabox}
\[
P(w_1, \dots, w_n) \approx \prod_{i=1}^{n} P(w_i \mid w_{i-1})
\]
\end{formulabox}
Ogni parola dipende solo dalla precedente (\emph{finestra di memoria limitata}).

\subsection{Stima MLE (Maximum Likelihood Estimation)}

\begin{formulabox}
\[
P_{\text{MLE}}(w_i \mid w_{i-1}) = \frac{C(w_{i-1}, w_i)}{C(w_{i-1})}
\]
\end{formulabox}

\textbf{Problema fondamentale}: se la coppia $(w_{i-1}, w_i)$ non è mai apparsa nel corpus di training, $P = 0$ e la probabilità dell'intera frase collassa a zero.

\subsection{Laplace Smoothing (Add-1)}

\begin{formulabox}
\[
P_{\text{Lap}}(w_i \mid w_{i-1}) = \frac{C(w_{i-1}, w_i) + 1}{C(w_{i-1}) + |V|}
\]
\end{formulabox}

dove $|V|$ è la dimensione del vocabolario. Nessuna probabilità zero, ma ridistribuisce eccessivamente la massa di probabilità dagli eventi frequenti a quelli non visti.

\subsection{Interpolazione lineare}

Combina modelli di ordine diverso per maggiore robustezza:

\begin{formulabox}
\[
P_{\text{interp}}(w_i \mid w_{i-1}) = \lambda_1 P(w_i) + \lambda_2 P(w_i \mid w_{i-1}) + \lambda_3 P(w_i \mid w_{i-2}, w_{i-1})
\]
\end{formulabox}
con $\lambda_1 + \lambda_2 + \lambda_3 = 1$. I $\lambda$ si ottimizzano su un corpus di development.

\subsection{Kneser-Ney Smoothing}

Il metodo più efficace in pratica. Introduce il concetto di \textbf{frequenza di continuazione}: quante parole diverse seguono $w_{i-1}$? Questo penalizza le parole che, pur frequenti in assoluto, appaiono sempre nello stesso contesto.

\begin{formulabox}
\[
P_{\text{KN}}(w_i \mid w_{i-1}) \propto \max\!\bigl(C(w_{i-1}, w_i) - d,\, 0\bigr) + \lambda(w_{i-1}) \cdot P_{\text{KN}}(w_i)
\]
\end{formulabox}
dove $d \in (0,1)$ è un parametro di discounting e $P_{\text{KN}}(w_i)$ è la probabilità \emph{di continuazione} (proporzione di bigrammi distinti in cui $w_i$ appare come secondo elemento).

\subsection{Perplexity}

La \textbf{perplexity} misura quanto un modello è ``sorpreso'' da un testo di test: più bassa è, meglio il modello rappresenta i dati.

\begin{formulabox}
\[
\text{PP}(W) = P(w_1 w_2 \cdots w_n)^{-1/n} = \left(\prod_{i=1}^{n} \frac{1}{P(w_i \mid w_{i-1})}\right)^{1/n}
\]
\end{formulabox}

Equivalentemente, tramite cross-entropy:
\[
\text{PP}(W) = 2^{H(W)} \qquad \text{dove} \quad H(W) = -\frac{1}{n}\sum_{i=1}^n \log_2 P(w_i \mid w_{i-1})
\]

\textbf{Interpretazione}: media geometrica dei fattori di inversione di probabilità. Un modello perfetto darebbe PP=1. Valori tipici: 30--100 per buoni modelli su testo inglese.

\subsection{Log-Likelihood}

\begin{formulabox}
\[
\text{LL}(W) = \sum_{i=1}^{n} \log P(w_i \mid w_{i-1})
\]
\end{formulabox}

Sempre $\leq 0$. Più alto (meno negativo) = modello migliore. Più stabile della perplexity per la classificazione (evita divisioni per zero e overflow).

\subsection{Type-Token Ratio (TTR)}

\begin{formulabox}
\[
\text{TTR} = \frac{|\text{Types}|}{|\text{Tokens}|}
\]
\end{formulabox}

Misura la ricchezza lessicale. Alto TTR = vocabolario vario (slang, abbreviazioni come nei Twitter). Basso TTR = testo ripetitivo (stile narrativo formale).

% =====================================================================
\section{Lab 4 -- N-grammi: Twitter vs Testo Letterario}

\subsection{Obiettivo}

Confrontare le proprietà statistiche e generative di due corpora molto diversi: tweet (social media) e prosa letteraria (Emma di Jane Austen, 1816).

\subsection{Dataset e preprocessing}

\begin{itemize}[leftmargin=1.5em]
  \item \textbf{Twitter}: 10.000 tweet (5.000 positivi + 5.000 negativi) da \texttt{nltk.corpus.twitter\_samples}.
  \item \textbf{Letterario}: \texttt{austen-emma.txt} da \texttt{nltk.corpus.gutenberg}.
\end{itemize}

\textbf{Pulizia Twitter}: rimozione URL (\verb|http\S+|), mention (\verb|@\w+|), hashtag (\verb|#|), caratteri non alfabetici.\\
\textbf{Pulizia letterario}: rimozione punteggiatura, lowercase.

\subsection{Pipeline principale (codice)}

\begin{lstlisting}
# Modello MLE (n=2) per generazione testo
train, vocab = padded_everygram_pipeline(2, [tokens_tweet])
modello = MLE(2)
modello.fit(train, vocab)

# Modello Laplace per perplexity (gestisce OOV)
train_lp, vocab_lp = padded_everygram_pipeline(2, [tokens_tweet])
modello_lp = Laplace(2)
modello_lp.fit(train_lp, vocab_lp)

# Probabilita' condizionale
p = modello.score('love', ['i'])  # P(love | i)

# Generazione testo
testo = modello.generate(30, random_seed=42)

# Perplexity (normalizzata per vocabolario)
pp = modello_lp.perplexity(word_tokenize(frase))
pp_norm = pp / len(modello_lp.vocab)
\end{lstlisting}

\subsection{Classificatore con Log-Likelihood}

\begin{lstlisting}
def log_likelihood(modello, tokens):
    score = 0.0
    for i in range(1, len(tokens)):
        prob = modello.score(tokens[i], [tokens[i-1]])
        score += math.log(prob) if prob > 0 else math.log(1e-10)
    return score

# Regola di decisione: vince il modello con LL piu' alta
ll_tw   = log_likelihood(modello_tweet, tokens)
ll_lett = log_likelihood(modello_lett,  tokens)
verdetto = "TWITTER" if ll_tw > ll_lett else "LETTERARIO"
\end{lstlisting}

\subsection{Risultati e conclusioni}

\begin{center}
\begin{tabular}{lll}
\toprule
\textbf{Metrica} & \textbf{Twitter} & \textbf{Letterario} \\
\midrule
Token totali & $\sim$120.000 & $\sim$160.000 \\
Parole uniche (types) & $\sim$20.000 & $\sim$8.000 \\
TTR & $\sim$0.17 (alto) & $\sim$0.05 (basso) \\
Lunghezza media frase & $\sim$10 parole & $\sim$22 parole \\
\bottomrule
\end{tabular}
\end{center}

\textbf{Osservazioni chiave}:
\begin{itemize}[leftmargin=1.5em]
  \item Twitter ha TTR più alto: vocabolario vario (slang, abbreviazioni, neologismi).
  \item Le frasi letterarie sono più lunghe e strutturate sintatticamente.
  \item MLE con OOV dà perplexity $= \infty$: il Laplace smoothing è \emph{indispensabile}.
  \item Il classificatore \textbf{log-likelihood} è più stabile e accurato della perplexity normalizzata.
  \item Per $n=3$ il modello soffre di sparsità (trigrammi non visti): $n=2$ è il miglior trade-off.
\end{itemize}

% =====================================================================
\section{Vector Space Model (VSM)}

\subsection{Bag-of-Words (BoW)}

Ogni documento diventa un vettore di frequenze di parole, ignorando l'ordine:

\[
d = [c(\text{word}_1, d),\; c(\text{word}_2, d),\; \dots,\; c(\text{word}_V, d)]
\]

\textbf{Limite}: ignora l'ordine, peso uniforme a tutte le parole.

\subsection{TF-IDF}

\textbf{TF-IDF} bilancia la frequenza locale (TF) con la rarità globale (IDF).

\begin{formulabox}
\[
\text{TF}(t, d) = \frac{C(t,d)}{|d|}
\qquad
\text{IDF}(t) = \log\!\frac{N}{\text{DF}(t)}
\qquad
\text{TF-IDF}(t, d) = \text{TF}(t,d) \times \text{IDF}(t)
\]
\end{formulabox}

dove $N$ = numero totale di documenti, $\text{DF}(t)$ = numero di documenti che contengono $t$.

\textbf{Effetti}: termini comuni in molti documenti (``the'', ``is'') $\to$ IDF basso $\to$ TF-IDF basso. Termini rari e discriminativi $\to$ IDF alto $\to$ TF-IDF alto.

\subsection{Cosine Similarity}

\begin{formulabox}
\[
\cos(d_1, d_2) = \frac{d_1 \cdot d_2}{\|d_1\| \cdot \|d_2\|} = \frac{\sum_i d_1[i]\cdot d_2[i]}{\sqrt{\sum_i d_1[i]^2}\cdot\sqrt{\sum_i d_2[i]^2}}
\]
\end{formulabox}

Range: $[0, 1]$ per vettori non negativi. $0$ = vettori ortogonali (completamente diversi); $1$ = stessa direzione. La \emph{normalizzazione} rende la misura invariante alla lunghezza del documento.

\subsection{Rocchio's Algorithm (Relevance Feedback)}

Aggiorna il vettore di query avvicinandolo ai documenti rilevanti e allontanandolo da quelli non rilevanti:

\begin{formulabox}
\[
Q' = \alpha \cdot Q + \beta \cdot \frac{1}{|R|}\sum_{D \in R} D \;-\; \gamma \cdot \frac{1}{|NR|}\sum_{D \in NR} D
\]
\end{formulabox}

Valori tipici: $\alpha=1,\; \beta=0.75,\; \gamma=0.25$.

\subsection{Latent Dirichlet Allocation (LDA)}

LDA è un modello probabilistico generativo per l'\textbf{analisi tematica} dei documenti.

\begin{defbox}
\textbf{Ipotesi}: ogni documento è una miscela di topic; ogni topic è una distribuzione su parole.
\end{defbox}

\textbf{Processo generativo}:
\begin{enumerate}[leftmargin=1.5em, label=\arabic*.]
  \item Per ogni documento $d$: estrai la distribuzione sui topic $\theta_d \sim \text{Dirichlet}(\alpha)$.
  \item Per ogni posizione di parola nel documento:
    \begin{enumerate}[label=\alph*.]
      \item Estrai il topic $z \sim \text{Multinomial}(\theta_d)$.
      \item Estrai la parola $w \sim \text{Multinomial}(\beta_z)$.
    \end{enumerate}
\end{enumerate}

\textbf{Vantaggi}: interpretabile, non supervisionato, gestisce polisemia (una parola è assegnata a topic diversi in contesti diversi). \textbf{Limiti}: bag-of-words (ignora l'ordine), numero di topic fissato a priori, lento.

% =====================================================================
\section{Word2Vec: Rappresentazioni Distribuite}

\subsection{Ipotesi distribuzionale}

\begin{defbox}
\textbf{Firth (1957)}: ``You shall know a word by the company it keeps.'' \\
Parole che appaiono in contesti simili hanno significati correlati.
\end{defbox}

\subsection{Architetture Word2Vec}

\subsubsection{CBOW (Continuous Bag of Words)}

\textbf{Input}: le parole di contesto (finestra di $\pm w$ parole). \textbf{Output}: la parola target centrale.

\begin{enumerate}[leftmargin=1.5em, label=\arabic*.]
  \item One-hot encode le parole di contesto.
  \item Embedding layer: lookup + media degli embedding.
  \item Output layer: softmax sul vocabolario.
  \item Loss: cross-entropy con la parola target.
\end{enumerate}

\textbf{Vantaggio}: veloce, efficiente su grandi corpora.

\subsubsection{Skip-gram}

\textbf{Input}: la parola target. \textbf{Output}: le parole di contesto. Inverso del CBOW.

\begin{enumerate}[leftmargin=1.5em, label=\arabic*.]
  \item One-hot encode la parola target.
  \item Embedding layer: lookup embedding target.
  \item Per ogni parola di contesto: output layer softmax.
  \item Loss: somma delle cross-entropy per ogni parola di contesto.
\end{enumerate}

\textbf{Vantaggio}: funziona meglio per parole rare; migliore qualità degli embedding su corpora piccoli.

\subsection{Negative Sampling}

\textbf{Problema}: il softmax sull'intero vocabolario ($O(|V|)$) è troppo costoso.

\textbf{Soluzione}: per ogni campione positivo $(w_{\text{target}}, w_{\text{contesto}})$, campionare $k$ parole \emph{negative} casuali e ottimizzare:

\begin{formulabox}
\[
\mathcal{L} = \log \sigma(u_t \cdot v_c) + \sum_{i=1}^{k} \log \sigma(-u_t \cdot v_{n_i})
\]
\end{formulabox}

dove $\sigma$ è la funzione sigmoide. Riduce la complessità da $O(|V|)$ a $O(k)$. Tipicamente $k = 5$--$20$.

\subsection{Proprietà degli embedding: word arithmetic}

\begin{formulabox}
\[
\vec{v}(\text{king}) - \vec{v}(\text{man}) + \vec{v}(\text{woman}) \approx \vec{v}(\text{queen})
\]
\[
\vec{v}(\text{Paris}) - \vec{v}(\text{France}) + \vec{v}(\text{Germany}) \approx \vec{v}(\text{Berlin})
\]
\end{formulabox}

Dimostra che lo spazio degli embedding codifica relazioni semantiche e sintattiche come trasformazioni lineari.

\subsection{Parametri di training}

\begin{center}
\begin{tabular}{llll}
\toprule
\textbf{Parametro} & \textbf{Descrizione} & \textbf{Valore tipico} & \textbf{Codice} \\
\midrule
\texttt{vector\_size} & Dimensione embedding & 100--300 & \texttt{vector\_size=100} \\
\texttt{window} & Finestra di contesto ($\pm w$) & 5--10 & \texttt{window=5} \\
\texttt{min\_count} & Frequenza minima & 2--5 & \texttt{min\_count=3} \\
\texttt{epochs} & Passi sul corpus & 5--20 & \texttt{epochs=20} \\
\texttt{sg} & 0=CBOW, 1=Skip-gram & -- & \texttt{sg=1} \\
\texttt{negative} & Negative samples & 5--20 & \texttt{negative=5} \\
\bottomrule
\end{tabular}
\end{center}

\subsection{Implementazione Gensim}

\begin{lstlisting}
from gensim.models import Word2Vec

model = Word2Vec(
    sentences=testi_tokenizzati,  # lista di liste di token
    vector_size=100, window=5, min_count=3,
    epochs=20, sg=1, seed=42
)

# Recupero embedding
vec = model.wv['love']  # array di 100 float

# Similarita' coseno tra due parole
sim = model.wv.similarity('love', 'heart')  # float in [-1,1]

# Parole piu' simili
simili = model.wv.most_similar('king', topn=5)
# [('queen', 0.87), ('prince', 0.82), ...]
\end{lstlisting}

\subsection{Analisi della deriva semantica (Semantic Shift)}

\textbf{Metrica di overlap} tra i $k$ vicini più simili in due ere diverse:

\begin{formulabox}
\[
\text{overlap}(w, \text{era}_1, \text{era}_2) = \frac{|\text{vicini}_{\text{era}_1}(w) \cap \text{vicini}_{\text{era}_2}(w)|}{k}
\]
\end{formulabox}

Interpretazione: $\text{overlap} \approx 1$ = significato stabile; $\text{overlap} \approx 0$ = deriva semantica completa.

\textbf{Delta similarità} tra coppie di parole:
\[
\Delta(w_1, w_2) = |\text{sim}_{\text{era}_1}(w_1,w_2) - \text{sim}_{\text{era}_2}(w_1,w_2)|
\]

\subsection{Limiti di Word2Vec}

\begin{itemize}[leftmargin=1.5em]
  \item \textbf{Embedding statici}: ogni parola ha un unico vettore, indipendentemente dal contesto (non gestisce la polisemia).
  \item \textbf{Finestra limitata}: considera solo il contesto locale ($\pm w$ parole).
  \item \textbf{Nessuna struttura sub-word}: ``walking'' e ``walk'' sono appresi indipendentemente.
  \item \textbf{Bias}: apprende pregiudizi presenti nel corpus di addestramento.
\end{itemize}

% =====================================================================
\section{Lab 8 -- Word2Vec su Testi Musicali (70s vs 2000s)}

\subsection{Obiettivo}

Analizzare il cambiamento semantico del vocabolario musicale tra gli anni '70 e gli anni 2000 tramite due modelli Word2Vec separati.

\subsection{Dataset e artisti}

\begin{center}
\begin{tabular}{ll}
\toprule
\textbf{Era} & \textbf{Artisti} \\
\midrule
Anni '70 & Led Zeppelin, Pink Floyd, Queen, Eagles, Fleetwood Mac, \ldots \\
Anni 2000 & Eminem, Coldplay, Linkin Park, Taylor Swift, Lady Gaga, \ldots \\
\bottomrule
\end{tabular}
\end{center}

Dataset: \texttt{spotify\_millsongdata.csv}.

\subsection{Pipeline di preprocessing}

\begin{lstlisting}
def pulisci_testo(testo):
    testo = re.sub(r'\[.*?\]', '', testo)   # rimuove [Chorus], [Verse]
    testo = re.sub(r'[^a-zA-Z\s]', '', testo)  # solo lettere
    testo = testo.lower()
    tokens = testo.split()
    # rimuove stopword e token corti (< 3 caratteri)
    return [t for t in tokens if t not in stop and len(t) >= 3]
\end{lstlisting}

\subsection{Addestramento e analisi}

\begin{lstlisting}
params = dict(vector_size=100, window=5, min_count=3,
              workers=4, epochs=20, seed=42)
model_70s   = Word2Vec(df_70s['tokens'].tolist(),   **params)
model_2000s = Word2Vec(df_2000s['tokens'].tolist(), **params)

# Parole simili per era
for parola in ['love', 'night', 'heart']:
    sim70  = [w for w,_ in model_70s.wv.most_similar(parola, topn=3)]
    sim00  = [w for w,_ in model_2000s.wv.most_similar(parola, topn=3)]
    print(f"{parola}: 70s={sim70}  2000s={sim00}")

# Analisi deriva: overlap dei vicini
for parola in parole_comuni:
    vicini_70  = set(w for w,_ in model_70s.wv.most_similar(parola, topn=10))
    vicini_00  = set(w for w,_ in model_2000s.wv.most_similar(parola, topn=10))
    overlap     = len(vicini_70 & vicini_00) / 10
    cambiamento = 1 - overlap
\end{lstlisting}

\subsection{Clustering K-Means + t-SNE}

\begin{lstlisting}
from sklearn.cluster import KMeans
from sklearn.manifold import TSNE

# Prendi top-200 parole per frequenza
parole  = model_70s.wv.index_to_key[:200]
vettori = np.array([model_70s.wv[w] for w in parole])

# K-Means con k=6 cluster
km = KMeans(n_clusters=6, random_state=42, n_init=10)
labels = km.fit_predict(vettori)

# t-SNE per visualizzazione 2D (100D -> 2D)
tsne = TSNE(n_components=2, random_state=42, perplexity=30)
coords_2d = tsne.fit_transform(vettori)
\end{lstlisting}

\subsection{Risultati e conclusioni}

\begin{itemize}[leftmargin=1.5em]
  \item \textbf{``love''}: negli anni '70 contesto romantico/metaforico; negli anni 2000 più diretto e personale.
  \item \textbf{``night''}: negli anni '70 mistero e oscurità; negli anni 2000 vita notturna urbana.
  \item Le parole caratteristiche degli anni '70 riflettono il rock psichedelico; quelle degli anni 2000 il rap e la musica elettronica.
  \item I cluster K-Means mostrano temi coerenti con i generi musicali delle rispettive ere.
  \item La \textbf{heatmap di similarità} evidenzia come le relazioni tra parole chiave (love, heart, god, soul\ldots) cambino tra le due ere.
  \item Metrica di overlap: una parola con overlap $< 50\%$ ha subito una deriva semantica significativa.
\end{itemize}

% =====================================================================
\section{Transformers e Rappresentazioni Contestuali}

\subsection{Limiti di Word2Vec che motivano i Transformers}

Word2Vec assegna a ogni parola \emph{un unico vettore}, indipendentemente dal contesto. Non gestisce la polisemia:

\begin{itemize}[leftmargin=1.5em]
  \item ``I went to the \textbf{bank} to withdraw cash.'' (banca finanziaria)
  \item ``I sat on the \textbf{bank} of the river.'' (riva del fiume)
\end{itemize}

Entrambe le occorrenze di ``bank'' ricevono lo stesso vettore in Word2Vec.

\subsection{Meccanismo di Attention}

\begin{defbox}
L'\textbf{attention} permette di pesare diversamente le parti dell'input in base al loro ruolo nel calcolo del significato di una specifica parola target.
\end{defbox}

\subsubsection{Scaled Dot-Product Attention}

\begin{formulabox}
\[
\text{Attention}(Q, K, V) = \text{softmax}\!\left(\frac{QK^\top}{\sqrt{d_k}}\right) V
\]
\end{formulabox}

Dove:
\begin{itemize}[leftmargin=1.5em]
  \item $Q$ (Query): cosa stiamo cercando?
  \item $K$ (Key): cosa offre ogni posizione?
  \item $V$ (Value): l'informazione da recuperare.
  \item $\sqrt{d_k}$: fattore di scala (evita gradienti saturi nel softmax).
\end{itemize}

\textbf{Intuitivamente}: per la parola ``bank'' in una frase finanziaria, il softmax darà pesi alti a parole come ``withdraw'', ``deposit'', ``money''; in un contesto fluviale darà pesi alti a ``river'', ``water'', ``sit''.

\subsubsection{Multi-Head Attention}

Esegue $h$ meccanismi di attention in parallelo con proiezioni apprese diverse:

\begin{formulabox}
\[
\text{MultiHead}(Q, K, V) = \text{Concat}(\text{head}_1, \dots, \text{head}_h) W^O
\]
\[
\text{head}_i = \text{Attention}(Q W^Q_i,\; K W^K_i,\; V W^V_i)
\]
\end{formulabox}

\textbf{Vantaggio}: ogni ``testa'' impara relazioni diverse (sintattiche, semantiche, coreference, \ldots).

\subsection{BERT (Bidirectional Encoder Representations from Transformers)}

\begin{defbox}
BERT è un Transformer \textbf{bidirezionale}: vede tutto il contesto (sia a sinistra che a destra) durante il pre-training.
\end{defbox}

\textbf{Pre-training task}:

\begin{enumerate}[leftmargin=1.5em, label=\arabic*.]
  \item \textbf{Masked Language Modeling (MLM)}: il 15\% dei token è mascherato; il modello deve predire il token originale.
  \begin{itemize}[leftmargin=1.5em, nosep]
    \item 80\% sostituito con \texttt{[MASK]}
    \item 10\% sostituito con parola casuale
    \item 10\% lasciato invariato
  \end{itemize}
  \item \textbf{Next Sentence Prediction (NSP)}: dato \texttt{[CLS] A [SEP] B [SEP]}, predire se B segue A nel testo.
\end{enumerate}

\textbf{Uso in downstream tasks}:
\begin{enumerate}[leftmargin=1.5em, label=\arabic*., nosep]
  \item Aggiungere token speciali: \texttt{[CLS]}, \texttt{[SEP]}.
  \item Passare attraverso BERT: embedding contestualizzati per ogni token.
  \item Il vettore \texttt{[CLS]} è usato come rappresentazione dell'intera frase.
  \item Fine-tuning sul task specifico.
\end{enumerate}

\subsection{GPT (Generative Pre-trained Transformer)}

\begin{defbox}
GPT è un Transformer \textbf{unidirezionale} (causale): può vedere solo i token precedenti, non quelli successivi.
\end{defbox}

\textbf{Pre-training}: language modeling autoregressivo:
\[
\mathcal{L} = -\sum_{t} \log P(w_t \mid w_1, \dots, w_{t-1})
\]

\subsection{BERT vs GPT: confronto}

\begin{center}
\begin{tabular}{lll}
\toprule
\textbf{Aspetto} & \textbf{BERT} & \textbf{GPT} \\
\midrule
Direzione & Bidirezionale & Unidirezionale (sinistra$\to$destra) \\
Pre-training & MLM + NSP & Language Modeling autoregressivo \\
Punti di forza & Comprensione, classificazione & Generazione testo \\
Embedding & Contestuali bidirezionali & Contestuali causali \\
\bottomrule
\end{tabular}
\end{center}

\subsection{Embedding contestuali vs statici}

Con BERT, ogni occorrenza di una parola riceve un vettore \emph{diverso} in base al contesto circostante:

\begin{lstlisting}
# Le due occorrenze di "bank" avranno embedding diversi
text1 = "I went to the bank to withdraw cash"
text2 = "I sat on the bank of the river"
# embedding1(bank) != embedding2(bank)  <- gestisce polisemia!
\end{lstlisting}

% =====================================================================
\section{Riepilogo: Confronti tra Approcci}

\subsection{Stemming vs Lemmatizzazione}

\begin{center}
\begin{tabular}{lll}
\toprule
& \textbf{Stemming} & \textbf{Lemmatizzazione} \\
\midrule
Tool & Regole morfologiche & Morfologia + dizionario \\
Input & Parola & Parola + PoS \\
Output & Radice (può non essere valida) & Lemma (parola valida) \\
Precisione & $\sim$70\% & $\sim$95\% \\
Velocità & Molto veloce & Lento \\
\bottomrule
\end{tabular}
\end{center}

\subsection{MLE vs Laplace vs Kneser-Ney}

\begin{center}
\begin{tabular}{llll}
\toprule
& \textbf{MLE} & \textbf{Laplace} & \textbf{Kneser-Ney} \\
\midrule
OOV & $P=0$ (fallisce) & $P>0$ & $P>0$ \\
Over-smoothing & No & Sì (eccessivo) & No \\
Complessità & Bassa & Bassa & Alta \\
Uso & Generazione & Perplexity, classif. & LM di produzione \\
\bottomrule
\end{tabular}
\end{center}

\subsection{WordNet vs Word2Vec vs BERT}

\begin{center}
\begin{tabular}{llll}
\toprule
\textbf{Aspetto} & \textbf{WordNet} & \textbf{Word2Vec} & \textbf{BERT} \\
\midrule
Costruzione & Manuale & Automatica (corpus) & Automatica (corpus) \\
Polisemia & Esplicita (synset) & Non gestita & Gestita (contestuale) \\
Similarità & Graph-based & Coseno & Coseno contestuale \\
Relazioni & Esplicite (is-a, \ldots) & Implicite & Implicite \\
Copertura & Limitata ($\sim$155k) & Dipende dal corpus & Vasta \\
Aggiornamento & Curato & Data-driven & Data-driven \\
\bottomrule
\end{tabular}
\end{center}

\subsection{N-grammi vs Word2Vec vs Transformers}

\begin{center}
\begin{tabular}{llll}
\toprule
\textbf{Aspetto} & \textbf{N-grammi} & \textbf{Word2Vec} & \textbf{Transformer} \\
\midrule
Contesto & Limitato ($n$ parole) & Finestra fissa & Intera sequenza \\
Polisemia & No & No (statico) & Sì (contestuale) \\
Training & Rapido & Moderato & Lento \\
Parametri & Pochi & $|V| \times d$ & Miliardi \\
Interpretabilità & Alta & Bassa & Molto bassa \\
Downstream NLP & Limitato & Buono & Eccellente \\
\bottomrule
\end{tabular}
\end{center}

% =====================================================================
\section{Formule di Riferimento}

\begin{center}
\begin{tabular}{ll}
\toprule
\textbf{Formula} & \textbf{Nome} \\
\midrule
$P_\text{MLE}(w|ctx) = C(ctx,w)/C(ctx)$ & MLE bigrammi \\
$P_\text{Lap}(w|ctx) = (C(ctx,w)+1)/(C(ctx)+|V|)$ & Laplace smoothing \\
$\text{PP}(W) = (\prod_i 1/P(w_i|w_{i-1}))^{1/n}$ & Perplexity \\
$\text{LL}(W) = \sum_i \log P(w_i|w_{i-1})$ & Log-Likelihood \\
$\text{TTR} = |\text{Types}|/|\text{Tokens}|$ & Type-Token Ratio \\
\midrule
$\text{TF}(t,d) = C(t,d)/|d|$ & Term Frequency \\
$\text{IDF}(t) = \log(N/\text{DF}(t))$ & Inv. Document Freq. \\
$\text{cos}(d_1,d_2) = (d_1 \cdot d_2)/(\|d_1\|\|d_2\|)$ & Cosine Similarity \\
\midrule
$\text{sim}_\text{WUP}(w_1,w_2) = 2\,\text{dep}(\text{LCA})/(\text{dep}(w_1)+\text{dep}(w_2))$ & Wu-Palmer \\
$\text{IC}(c) = -\log P(c)$ & Information Content \\
$\text{sim}_\text{Lin}(w_1,w_2) = 2\,\text{IC}(\text{LCA})/(\text{IC}(w_1)+\text{IC}(w_2))$ & Lin \\
\midrule
$\text{Attention}(Q,K,V) = \text{softmax}(QK^\top/\sqrt{d_k})\,V$ & Scaled Dot-Product \\
$\mathcal{L}_\text{neg} = \log\sigma(u_t\cdot v_c)+\sum_i\log\sigma(-u_t\cdot v_{n_i})$ & Negative Sampling \\
\bottomrule
\end{tabular}
\end{center}

% =====================================================================
\vfill
\begin{center}
\color{darkgray}\small\itshape
Documento preparato a scopo di studio per l'esame di TLN -- A.A. 2025/26\\
Università degli Studi di Torino -- Parte 2 (Prof.\ Radicioni)
\end{center}

\end{document}
